\documentclass[10pt,a4paper]{amsart}
\usepackage[margin=19mm]{geometry}
\usepackage{amsmath,amssymb,mathtools}
\usepackage[scr=rsfs,cal=euler]{mathalfa}
\usepackage{xfrac}
\usepackage[unicode,hidelinks,bookmarks=false]{hyperref}
\newcommand{\ie}{i.e.\ }
\newcommand{\cf}{cf.\ }
\newcommand{\resp}{respectively\ }
\newcommand{\Subsection}{\subsection}
\providecommand{\nobreakdash}{\nobreak\hbox{-}}
\setlength{\emergencystretch}{2em}
\multlinegap0pt
\usepackage{amssymb}
\newtheorem{theorem}{Theorem}
\newcommand{\Int}[1]{$#1$\textnormal{-interval-PCG}}
\title{All Graphs with at Most 8 Nodes are \Int{2}s}
\author{Tiziana Calamoneri, Angelo Monti, Fabrizio Petroni}
\date{Source: arXiv:2202.13844v4 (CC BY 4.0)}
\begin{document}
\maketitle

\noindent\emph{Source and licence.} All Graphs with at Most 8 Nodes are \Int{2}s. Version arXiv:2202.13844v4 (CC BY 4.0), distributed by arXiv under the Creative Commons Attribution 4.0 International licence, http://creativecommons.org/licenses/by/4.0/, as recorded in the arXiv licence metadata for this exact version and shown on the abstract page https://arxiv.org/abs/2202.13844v4. Full licence notice: \texttt{licenses/arxiv-2202.13844-CC-BY-4.0.txt}.\par\smallskip

\noindent\textit{Editorial scope.} Counted unit: the article's theorem theo:almost (source0.tex lines 170--173). The article's complete original proof of it is delivered with it (source0.tex lines 174--197, one \texttt{proof} environment of 24 lines). No mathematics is written, repaired or re-derived: the statement and the proof are the article's own lines, in the article's own notation, and no retained line is substituted or re-flowed.  No further source-local context is needed: every cross-reference of every retained line resolves inside the two delivered spans.  Nothing was omitted from any retained span, so the package carries no omission record.  No retained line contains a citation, so the wrapper carries no bibliography and no bibliography entry.  No retained line contains a figure environment, an image-inclusion command or drawing code, so no image is delivered and the package declares no asset dependency.  Editorial formatting only, in addition: an AMS \texttt{amsart} preamble that restates the article's own declarations and macros used by the retained lines, namely the environments \texttt{theorem} on separate counters, together with the standard packages those lines need.  The retained environments print in this document as Theorem 1 in order of appearance, whereas the article prints the same items with the numbering of its own full document.  The complete native source of the article as distributed in the arXiv e-print for this exact version is bundled unchanged as \texttt{sources/122640/source0.tex}.
\begin{theorem}
\label{theo:almost}
Let $G$ be a graph with an almost universal node $u$. If the subgraph $G'$ of $G$, obtained from $G$ by removing $u$, is  an \Int{k}, then $G$ is a \Int{(k+1)}.
\end{theorem}

\begin{proof}
Let $G'=PCG(T,I_1,\ldots I_k)$ with $I_i=[a_i, b_i]$, $1\leq i\leq k$.
%Using property \label{prop:interi} we can assume that $T$ is weighted  by natural numbers.
Starting from tree $T$ we now create a new tree $T_1$ by adding a weight $2$
 to all the edges of $T$ incident to the leaves.
In this way, the distances between any two leaves in $T$ are augmented by exactly $4$ in $T_1$ and hence
$G'=PCG(T_1, I'_1,\ldots I'_k)$ where $I_i=[a_i+4,b_i+4]$ and the weight of each edge incident to a leaf is strictly larger than 1.

Let $v$ be the only node not connected with
the almost universal node $u$ of $G$, and let $x$ be the parent of leaf $v$ in $T_1$; call  $c$ the weight of edge $(v,x)$ in $T_1$. (Note that $c \geq 2$).
Now construct a new tree $T_2$ obtained from $T_1$ deleting edge $(v,x)$ and adding a  dummy internal node $y$ connected with $v$ and $x$, where edge $(v,y)$ has weight $1$ and edge $(y,x)$ has weight $c-1$. Obviously, again $G'=PCG(T_2,I'_1,\ldots,I'_k)$.


%We finally construct a new tree $T_3$ obtained by $T_2$ adding to it node $u$ as leaf and define a new interval $I_2$.
%
Let $p$ be the maximum distance between the leaves of $T_2$ and hence it is not restrictive to assume that for the right extreme of $I_k$ is $b_k+4 \leq p$.
Finally construct $T_3$ adding to $T_2$ new leaf $u$ connected to $y$ with an edge of weight $p$.
Note that $u$ has distance $p+1$ from $v$ in $T_3$.
Moreover, the distance in $T_3$ between $u$ and any other leaf is less than $2p$ and at least  $p+2$.
Indeed, the weight of edge $(u,y)$ is $p$ and the weight of the unique incident edge of any other leaf is at least 2.
So, we set $I'_{k+1}=[p+2,2p]$.

Easily $I'_1, \ldots I'_k$ and $I_{k+1}$ are disjoint and $G=(T_3,I'_1,\ldots,I'_k,I'_{k+1})$ and hence it is a \Int{(k+1)}.
\end{proof}

\end{document}
